\documentclass[10pt]{article}
\usepackage[margin=20mm]{geometry}
\usepackage{amsmath,amssymb,mathtools,fontspec,unicode-math,xurl,hyperref}
\IfFontExistsTF{DejaVu Serif}{\setmainfont{DejaVu Serif}}{\setmainfont{Latin Modern Roman}}
\setmathfont{latinmodern-math.otf}
\setlength{\parindent}{0pt}\setlength{\parskip}{6pt}\sloppy
\title{Scientific initial review of M25-0025}
\author{}
\date{25 October 2026}
\begin{document}
\maketitle
Abstract: This review of Kisin and Zhou’s arXiv:2103.09945v2 finds no substantive issue in the inspected conjugacy-class descent, residue-field preservation, and selected explicit curve calculation. The rational point of the semisimple conjugacy quotient must be distinguished from a rational element of the group. A cited companion theorem’s hypotheses were checked against its primary source, and the ramified \(SU_3\) matrix calculation was independently verified exactly. The full deformation and integral-model machinery is not certified.\par
\section*{Paper identity and coverage}
The paper is Independence of \(\ell\) for Frobenius conjugacy classes attached to abelian varieties, by Mark Kisin and Rong Zhou, DOI https://doi.org/10.4007/annals.2025.202.3.3. The inspected source is the manifest-pinned author manuscript arXiv:2103.09945v2, 7 October 2024, 63 PDF pages. Page numbers below are physical PDF page numbers starting at one. All eight input hashes matched the manifest. Journal-final equivalence has not been established.\par
Text was read on pages 1–6, 33–34, and 42–61. Images actually viewed were pages 1, 2, 34, 43, 46, 49, 52, 54, 55, 58, 59, and 60. Source lines 1950–1995 in `l-independence\_\_arXiv.tex` were also read. This is a systematic review of selected claims and their dependencies, not a complete line-by-line verification of Sections 2–4. No author code was executed.\par
\section*{Assumptions and exact theorem scope}
Let \(A/E\) be an abelian variety over a number field with a fixed complex embedding, and let \(G/\mathbb Q\) be its Mumford–Tate group. The base field is first enlarged so that the Galois representations factor through \(G(\mathbb Q_\ell)\) for every \(\ell\). For a place \(v\mid p\) of good reduction with \(p>2\), Theorems 1.1 and 6.2.7 produce \[\gamma\in\operatorname{Conj}_G(\mathbb Q),\qquad
 \gamma_{\mathbb Q_\ell}
   =\chi_G\bigl(\rho_\ell^G(\operatorname{Frob}_v)\bigr)
       \quad(\ell\ne p).\] Here \(\operatorname{Conj}_G=\operatorname{Spec}\mathbb Q[G]^G\) records geometric semisimple conjugacy classes. A rational point of this quotient need not have a representative in \(G(\mathbb Q)\), nor does it by itself identify rational conjugacy classes in \(G(\mathbb Q_\ell)\).\par
The stronger rational-representative statement on page 60 assumes that \(G^{\rm der}\) is simply connected and \(G_{\mathbb Q_p}\) is quasi-split, and permits one exceptional prime away from \(p\). The main theorem makes neither the \(p=2\) assertion nor the \(\ell=p\) crystalline assertion. Good reduction is essential to the stated theorem. These are boundaries of the result, not defects.\par
The intermediate Shimura statement assumes a strongly acceptable triple (Definition 4.2.2): Hodge type, the indicated Weil restrictions of split derived groups, an induced-torus condition on the specified central subgroup, and a very special parahoric. The \(\mu\)-ordinary lifting result is used under these hypotheses, not asserted here for every Shimura variety or every level.\par
\section*{Checked descent and compatibility arguments}
\subsection*{Why curve points must preserve residue fields}
Proposition 5.2.3 and Corollary 5.2.7 construct a smooth curve over \(\mathbb F_q\) with an \(\mathbb F_q\)-rational point above the chosen \(\mathbb F_q\)-point of the stratum. If the point above it had residue field \(\mathbb F_{q^d}\), its Frobenius would give the \(d\)-th power of the original element. Equality of powers would not in general recover the original semisimple class. The actual statements retain \(d=1\); the introduction explicitly identifies this issue on page 4. Thus the proof’s use of curves is not interchangeable with an arbitrary curve after a field extension.\par
The later finite extension of the number field in the proof of Theorem 6.2.7, page 59, also preserves the residue field at the selected place. This can be justified directly. The inverse image \(U\) of the chosen compact open subgroup under the continuous Galois representation is an open subgroup containing the chosen lift \(\widetilde\sigma_q\) of geometric Frobenius. Take the finite fixed field of \(U\). The local decomposition subgroup intersected with \(U\) still contains that lift, so its image in the procyclic residue Galois group is all of that group. The chosen place therefore has residue degree one. An arbitrary finite extension would not have this property; the specific subgroup construction does.\par
\subsection*{Companions and semisimplification}
Proposition 5.3.5, page 54, uses Chin’s Theorem 4.6. The primary text of that theorem was consulted: it applies to semisimple lisse sheaves on smooth curves whose algebraic Frobenius eigenvalues are units at every finite place away from \(p\). It allows an enlargement of the coefficient number field. Thus being a unit only at the initial \(\ell\)-adic place, in isolation, would not state the complete hypothesis.\par
In the present geometric setting the stronger property follows. For the faithful abelian-variety representation the eigenvalues, or their inverses according to the homology convention, are units away from \(p\): the Weil eigenvalues and \(q/\alpha\) are algebraic integers. At every prime not dividing \(q\), both valuations are nonnegative and sum to zero. Tensor products, duals, and subquotients preserve this property. Every algebraic representation of \(G\) occurs in the tensor category generated by a faithful representation. Consequently the representations used in Proposition 5.3.5 have the required property.\par
One may apply the companion theorem to semisimplifications, since characteristic polynomials of every group element are unchanged. Equality on the dense open part of the curve then identifies the semisimplifications by Chebotarev, giving equality at the selected point. This is sufficient for the semisimple conjugacy quotient; no assertion about unipotent parts is required. No missing finite-order determinant hypothesis is imposed by the cited version of the companion theorem.\par
\subsection*{Descent of the common algebraic class to \(\mathbb Q\)}
The final step on page 55 uses more than rationality in a single completion. Let \(F\) be the residue field of the common algebraic point, \(F'/\mathbb Q\) its Galois closure, and put \(\Gamma=\operatorname{Gal}(F'/\mathbb Q)\), \(H=\operatorname{Gal}(F'/F)\). An embedding \(F\hookrightarrow\mathbb Q_\ell\), for every unramified \(\ell\ne p\), implies that its Frobenius conjugacy class in \(\Gamma\) meets \(H\). Chebotarev therefore forces every conjugacy class to meet \(H\). Omitting the one prime \(p\) and finitely many ramified primes does not affect this inference.\par
For clarity, the finite-group conclusion has a short complete proof. The transitive action of \(\Gamma\) on \(\Gamma/H\) has average number of fixed points equal to one, by orbit counting. If \(H\) meets every conjugacy class, every element fixes at least one coset. If \([\Gamma:H]>1\), the identity fixes more than one, contradicting that average. Hence \(H=\Gamma\), \(F=\mathbb Q\), and the class is rational. No issue was found in this descent.\par
\subsection*{Returning from the auxiliary group}
Lemma 6.2.6, pages 58–59, uses \[G\hookrightarrow H\subset\operatorname{Res}_{F/\mathbb Q}G_F.\] Over \(\overline{\mathbb Q}\), the first embedding is diagonal. If two diagonal elements become conjugate in the product group, any one component of the conjugating element already conjugates them in \(G\). The subsequent central modification \(H_1=H\times^{Z_H}T\) also preserves geometric conjugacy classes of \(H\): conjugation of \((h,1)\) by a pair \((h',t)\) changes only the first component by \(h'\). Equality modulo the central anti-diagonal does not introduce an extra ambiguity, since equality of the second components forces that central element to be the identity. Thus the asserted injection of geometric conjugacy classes is valid. Since it is Galois-equivariant, a rational auxiliary class with a preimage has a unique Galois-fixed preimage. This verifies the final descent mechanism, conditional on construction of the auxiliary Shimura datum.\par
\section*{Exact check of the ramified \(SU_3\) calculation}
Page 52, Case (3) of Proposition 5.2.3, uses \(p>2\), a uniformizer \(u\) with \(\tau(u)=-u\), and a nonzero constant \(x\). The defining unipotent condition for \(u_-(x,-x^2/2)\) holds because \(x^2-x^2/2-x^2/2=0\). The three matrices in the displayed factorization multiply to \[R=
 \begin{pmatrix}
 0&0&-2/(ux^2)\\
 0&1&-2/x\\
 ux^2/2&ux&-u
 \end{pmatrix}.\] An independent symbolic calculation verified this equality, \(\det R=1\), and \(R^{\tau\,t}JR=J\), where \(J\) is the anti-diagonal identity matrix. For \(s=\operatorname{diag}(1,1,u)\), one obtains \[s^{-1}Rs=
 \begin{pmatrix}
 0&0&-2/x^2\\
 0&1&-2u/x\\
 x^2/2&x&-u
 \end{pmatrix}.\] All entries are integral in the \(u\)-valuation, since \(x\) and \(2\) are units. Its determinant is one, so its inverse is also integral. This establishes membership in the indicated non-standard parahoric for the displayed matrix. The formula for the curve itself extends at \(x=0\), where it is the identity. The denominators in the factorization are used only for \(x\ne0\).\par
This check verifies a genuine local ingredient. It does not verify all root-orbit and Schubert-order reductions in Proposition 5.2.3, the entire classification of parahorics, or every stack-theoretic passage in Corollary 5.2.7.\par
\section*{Dependence, outcome, and limits}
The proof proceeds from canonical liftings on the \(\mu\)-ordinary locus to compatible systems along curves, then to all strata and finally to the abelian-variety statement through auxiliary groups. The role of the \(\mu\)-ordinary lifting in Corollary 4.4.13 was traced; the full \(p\)-divisible-group deformation construction was not independently audited.\par
Check | Judgment | Remaining scope\par
Quotient versus group element | Scope clarified | Stronger refinements retain extra hypotheses.\par
Curve companion step | No Issue Found | External theorem statement checked; proof not reproduced.\par
Rational class descent | No Issue Found | Finite-group argument supplied.\par
Residue-field preservation | No Issue Found | Specific open subgroup required.\par
\(SU_3\) matrix | No Issue Found | Exact identity, determinant, form, integrality checked.\par
Full theorem | Not certified | Integral-model and deformation machinery remain.\par
Acceptance / adoption | Pending / not requested | No human or author decision asserted.\par
No substantive defect is established in these checks, and no Problem section is manufactured.\par
The arXiv metadata lists v2 as latest and states that some local-model results moved to arXiv:. That external construction was not audited in full. No additional correction notice was identified in the consulted primary metadata or targeted search; absence everywhere is not claimed. The final journal PDF was not compared.\par
\section*{Supporting checks}
It contains input validation, extracted pages, source-page PNGs, the actual inspection ledger, source fragments, `su3\_check.py`, `su3\_result.json`, execution evidence, and the primary-reference check. Python with `-B` executed the reviewer-authored script; process 23816 exited zero, using SymPy 1.13.3. The rendering process 50912 exited zero. No author code was run or imported. The exact matrix assertions are rational identities, not floating-point experiments.\par
All source and evidence hashes are bound in provenance. All scientific analysis, computations, static source inspection, and the listed page inspections were performed by the current assistant. Human acceptance remains pending.\par
The authorized followup uses the existing Tectonic 0.17.0 runtime. Exact-source build logs, input and PDF hashes, and report-page inspection records are retained in the accompanying provenance and followup evidence. No compiler was installed by this conversation. The source-PDF image checks listed above remain a separate record.\par
9 M. Kisin and R. Zhou, Independence of \(\ell\) for Frobenius conjugacy classes attached to abelian varieties, arXiv:2103.09945v2, 7 October 2024. https://arxiv.org/abs/2103.09945. C. Chin, Independence of \(\ell\) of monodromy groups, J. Amer. Math. Soc. 17 (2004), 723–747; arXiv:math/0206147v4, Section 4, especially Theorem 4.6. https://arxiv.org/html/math/0206147v4. Annals of Mathematics, metadata for DOI 10.4007/annals.2025.202.3.3, https://annals.math.princeton.edu/2025/202-3/p03, consulted 25 October 2026.\par
\end{document}
